\section{DoS Attacks}
\paragraph{Goal :} overload or crash a server to make the service unavailable to legitimate users. There can be \textbf{semantic} attacks ("smart") or \textbf{brute-force} attacks.

\paragraph{Semantic attacks } overload a server by sending knowledge-specific requests. These can include : 
\begin{itemize}
    \item computational intensive SQL queries
    \item error generating requests
    \item part of an HTTP request so that the server waits for the rest
\end{itemize}

\paragraph{Brute-force attacks } overload a server by sending many "random" requests. These can include : 
\begin{itemize}
    \item Exhausting number of TCP connections (SYN flooding)
    \item Exhausting the number of application sessions
\end{itemize}

\paragraph{Distributed DoS attack (DDoS)}  can be executed by combining many attacker sources with different IP addresses.


\paragraph{Reflected DoS attack (RDoS)}

A reflected DoS attack is performed by spoofing the victim's IP address and using it as the source address of requests sent to third-party servers (e.g., DNS servers). The servers then send their responses to the victim instead of the attacker.

\vspace{0.3em}

This technique provides two advantages:

\begin{itemize}
    \item \textbf{Reflection:} the attack traffic appears to originate from legitimate servers rather than from the attacker.
    \item \textbf{Amplification:} the response is often larger than the request. As a result, a small amount of traffic sent by the attacker can generate a much larger volume of traffic towards the victim.
\end{itemize}

\vspace{0.3em}

\paragraph{Network Ingress Filtering}

Network ingress filtering is a defense against such attacks based on IP spoofing. It consists of configuring firewalls to verify that the source IP address of a packet is consistent with the network from which the packet originates.

This filtering is typically performed by the attacker's ISP or edge router, as close as possible to the source of the traffic.